% -*- coding: utf-8 -*-
% XeLaTeX 编译
\documentclass[a3paper, landscape, 11pt, twocolumn]{article}

% 引入宏包
\usepackage{fontspec}
\usepackage[UTF8]{ctex}
\usepackage{geometry}
\usepackage{listings}
\usepackage{xcolor}
\usepackage{enumitem}
\usepackage{tabularx}
\usepackage{makecell}
\usepackage{fancyhdr}
\usepackage{amssymb}
\usepackage{lastpage}

% 设置页面边距 (A3横向双栏)
\geometry{left=2cm, right=2cm, top=2.5cm, bottom=2.5cm, columnsep=1.8cm}

% 颜色定义
\definecolor{codebg}{RGB}{245,245,245}
\definecolor{keywordcolor}{RGB}{0,0,150}
\definecolor{commentcolor}{RGB}{0,128,0}
\definecolor{stringcolor}{RGB}{163,21,21}

% Java代码块样式设置
\lstdefinestyle{JavaExam}{
    language=Java,
    basicstyle=\ttfamily\small,
    backgroundcolor=\color{codebg},
    keywordstyle=\color{keywordcolor}\bfseries,
    commentstyle=\color{commentcolor},
    stringstyle=\color{stringcolor},
    showstringspaces=false,
    breaklines=true,
    frame=single,
    rulecolor=\color{black!50},
    tabsize=4,
    xleftmargin=0.5em,
    xrightmargin=0.5em,
    aboveskip=0.5em,
    belowskip=0.5em,
    numbers=left,
    numberstyle=\tiny\color{gray}
}

% 列表样式
\newlist{options}{enumerate}{1}
\setlist[options]{label=\Alph*., nosep, leftmargin=*, itemsep=0.3em}

% 填空题横线
\newcommand{\blank}{\underline{\hspace{2cm}}}

% 页眉页脚设置
\pagestyle{fancy}
\fancyhf{}
\fancyhead[C]{\Large\bfseries 《Java程序设计》期末考试试卷}
\fancyfoot[C]{第 \thepage\ 页 共 \pageref{LastPage}\ 页}
\fancyfoot[R]{得分：\underline{\hspace{2cm}}}
\fancyfoot[L]{题号：\underline{\hspace{1cm}}}
\fancyheadoffset{0cm}
\fancyfootoffset{0cm}

\begin{document}

% 试卷头部
\twocolumn[{
			\begin{center}
				\vspace*{-1cm}
				{\huge\bfseries 《Java程序设计》期末考试试卷（D 卷）}\\[0.8em]
				{\large 考试时间：120 分钟 \hspace{2em} 总分：150 分}\\[1em]
				\hrule
				\vspace{0.8em}
				\begin{tabularx}{0.95\linewidth}{|c|X|c|X|c|X|}
					\hline
					\makecell{题号} & \makecell{一                                                                        \\(40分)} & \makecell{二\\(20分)} & \makecell{三\\(20分)} & \makecell{四\\(30分)} & \makecell{五\\(40分)} \\
					\hline
					得分            & \hspace{1.5cm} & \hspace{1.5cm} & \hspace{1.5cm} & \hspace{1.5cm} & \hspace{1.5cm} \\
					\hline
				\end{tabularx}\\[1em]
				{\large 姓名：\underline{\hspace{3cm}} \hspace{1.5em} 学号：\underline{\hspace{4cm}} \hspace{1.5em} 班级：\underline{\hspace{3cm}}}
				\vspace{0.8em}
				\hrule
				\vspace{1em}
			\end{center}
		}]

\section*{一、单项选择题（共 20 小题，每小题 2 分，共 40 分）}
\textit{（在每小题给出的四个选项中，只有一项是符合题目要求的）}

\begin{enumerate}[label=\arabic*., itemsep=0.8em]
	\item 下列关于 JDK、JRE、JVM 的关系，正确的是？
	      \begin{options}
		      \item JDK 包含 JRE，JRE 包含 JVM
		      \item JRE 包含 JDK，JDK 包含 JVM
		      \item JVM 包含 JRE，JRE 包含 JDK
		      \item 三者相互独立，无包含关系
	      \end{options}

	\item 在 Java 中，以下哪个选项可以正确定义一个名为 `MAX` 的常量？
	      \begin{options}
		      \item public int MAX = 100;
		      \item public static int MAX = 100;
		      \item public static final int MAX = 100;
		      \item final int MAX = 100;
	      \end{options}

	\item 关于方法重写（Override）的描述，正确的是？
	      \begin{options}
		      \item 重写方法的访问权限可以比父类更严格
		      \item 重写方法的返回值类型必须与父类完全相同
		      \item 重写方法不能抛出比父类更多的受检查异常
		      \item 私有方法可以被重写
	      \end{options}

	\item 下列哪个不是 Java 中访问控制修饰符？
	      \begin{options}
		      \item public
		      \item private
		      \item protected
		      \item internal
	      \end{options}

	\item 关于 StringBuilder 和 StringBuffer 的说法，正确的是？
	      \begin{options}
		      \item StringBuilder 是线程安全的
		      \item StringBuffer 是线程安全的
		      \item 两者性能相同
		      \item StringBuilder 的所有方法都使用 synchronized 修饰
	      \end{options}

	\item 下列哪个异常是受检查异常（checked exception）？
	      \begin{options}
		      \item NullPointerException
		      \item IndexOutOfBoundsException
		      \item IOException
		      \item ArithmeticException
	      \end{options}

	\item 关于 Java 中的 `final` 关键字，说法错误的是？
	      \begin{options}
		      \item final 修饰的类不能被继承
		      \item final 修饰的方法不能被重写
		      \item final 修饰的变量必须在声明时初始化
		      \item final 修饰的引用变量不能再指向其他对象
	      \end{options}

	\item 下列哪个类是所有 Java 集合框架的根接口？
	      \begin{options}
		      \item Collection
		      \item List
		      \item Set
		      \item Map
	      \end{options}

	\item 关于 Java 泛型，以下说法正确的是？
	      \begin{options}
		      \item 泛型类可以有多个类型参数
		      \item 泛型类型参数可以是基本数据类型
		      \item 泛型在运行时可以获取具体类型
		      \item 泛型数组可以直接创建，如 `new T[10]`
	      \end{options}

	\item 在 Java 线程中，下列哪个方法可以使当前线程让出 CPU，但不会释放锁？
	      \begin{options}
		      \item wait()
		      \item sleep()
		      \item yield()
		      \item join()
	      \end{options}

	\item 下列哪个接口用于实现对象的比较（自然排序）？
	      \begin{options}
		      \item Comparator
		      \item Comparable
		      \item Compare
		      \item Sortable
	      \end{options}

	\item 关于 Java I/O 中 `BufferedReader` 的说法，正确的是？
	      \begin{options}
		      \item 它是字节输入流
		      \item 它只能读取文本文件
		      \item 它可以提供缓冲功能，提高读取效率
		      \item 它的构造方法可以直接传入文件路径
	      \end{options}

	\item 下列哪个模式是 Java 事件处理机制采用的？
	      \begin{options}
		      \item 观察者模式
		      \item 工厂模式
		      \item 单例模式
		      \item 代理模式
	      \end{options}

	\item 关于 `java.lang.Object` 类的方法，正确的是？
	      \begin{options}
		      \item `clone()` 方法默认实现深克隆
		      \item `equals()` 方法默认比较对象引用
		      \item `finalize()` 方法一定会被调用
		      \item `notify()` 方法唤醒所有等待线程
	      \end{options}

	\item 在 Java 中，`transient` 关键字的作用是？
	      \begin{options}
		      \item 标记变量为不可序列化
		      \item 标记变量为临时变量，不会被垃圾回收
		      \item 标记变量为易失变量，保证可见性
		      \item 标记变量为线程局部变量
	      \end{options}

	\item 下列哪个类可以实现线程安全的 List？
	      \begin{options}
		      \item ArrayList
		      \item LinkedList
		      \item Vector
		      \item AbstractList
	      \end{options}

	\item 关于 Java 反射，以下说法错误的是？
	      \begin{options}
		      \item 可以在运行时获取类的构造方法
		      \item 可以在运行时调用私有方法
		      \item 反射会降低程序性能
		      \item 反射不能访问静态成员
	      \end{options}

	\item 在 JDBC 中，`ResultSet` 对象的默认游标位置是？
	      \begin{options}
		      \item 第一行之前
		      \item 第一行
		      \item 最后一行之后
		      \item 随机位置
	      \end{options}

	\item 关于 Swing 中 `JOptionPane` 的描述，正确的是？
	      \begin{options}
		      \item 只能显示消息对话框
		      \item 可以创建输入对话框
		      \item 所有方法都是非静态的
		      \item 必须作为 JFrame 的成员使用
	      \end{options}

	\item 在 TCP 网络编程中，`Socket` 类的 `getInputStream()` 方法返回的是？
	      \begin{options}
		      \item OutputStream
		      \item InputStream
		      \item byte 数组
		      \item 字符串
	      \end{options}
\end{enumerate}

\section*{二、判断题（共 10 小题，每小题 2 分，共 20 分）}
\textit{（正确的选 A，错误的选 B）}

\begin{enumerate}[label=\arabic*., itemsep=0.6em]
	\item Java 源文件名必须与 public 类名相同。 \underline{\hspace{1cm}}（A/B）
	\item 在同一个类中，方法重载的参数列表必须不同，但返回值类型可以相同也可以不同。 \underline{\hspace{1cm}}（A/B）
	\item 抽象类中的抽象方法必须被非抽象子类实现。 \underline{\hspace{1cm}}（A/B）
	\item 接口中的成员变量默认是 `public static final` 的。 \underline{\hspace{1cm}}（A/B）
	\item 使用 `throw` 关键字抛出异常时，必须使用 `throws` 在方法声明中声明。 \underline{\hspace{1cm}}（A/B）
	\item `HashSet` 允许存储重复元素，但存储顺序不确定。 \underline{\hspace{1cm}}（A/B）
	\item 泛型方法可以在非泛型类中定义。 \underline{\hspace{1cm}}（A/B）
	\item `Thread` 类的 `run()` 方法直接启动新线程。 \underline{\hspace{1cm}}（A/B）
	\item `File` 类的 `delete()` 方法只能删除文件，不能删除目录。 \underline{\hspace{1cm}}（A/B）
	\item 在 UDP 通信中，`DatagramSocket` 负责发送和接收数据报。 \underline{\hspace{1cm}}（A/B）
\end{enumerate}

\section*{三、填空题（共 5 小题，每空 2 分，共 20 分）}

\begin{enumerate}[label=\arabic*., itemsep=0.8em]
	\item Java 程序从源代码到运行的过程包括：编译（使用 \blank 命令）和运行（使用 \blank 命令）。
	\item Java 中的 8 种基本数据类型中，占用空间最小的是 \blank，占用空间最大的是 \blank。
	\item 面向对象设计原则中，单一职责原则要求一个类只负责 \blank，开闭原则要求对扩展 \blank，对修改关闭。
	\item 线程同步中，使用 \blank 关键字修饰方法或代码块，使用 \blank 类可以实现显式锁。
	\item Java 8 中，接口可以包含 \blank 方法和 \blank 方法，它们可以有方法体。
\end{enumerate}

\section*{四、程序运行结果题（共 5 小题，每小题 6 分，共 30 分）}
\textit{（阅读下列程序，写出运行后的输出结果）}

\begin{enumerate}[label=\arabic*., itemsep=1em]
	\item
	      \begin{lstlisting}[style=JavaExam]
public class TestA {
    public static void main(String[] args) {
        int[] arr = {3, 1, 4, 1, 5};
        int sum = 0;
        for (int n : arr) {
            if (n % 2 == 0) continue;
            sum += n;
        }
        System.out.println(sum);
    }
}
\end{lstlisting}
	      \textbf{输出：}\underline{\hspace{5cm}}

	\item
	      \begin{lstlisting}[style=JavaExam]
class Parent {
    String name = "Parent";
    public String getName() {
        return name;
    }
}
class Child extends Parent {
    String name = "Child";
    public String getName() {
        return name;
    }
}
public class TestB {
    public static void main(String[] args) {
        Parent p = new Child();
        System.out.println(p.name);
        System.out.println(p.getName());
    }
}
\end{lstlisting}
	      \textbf{输出：}\underline{\hspace{5cm}}

	\item
	      \begin{lstlisting}[style=JavaExam]
public class TestC {
    public static void main(String[] args) {
        try {
            System.out.print("A");
            throw new RuntimeException();
        } catch (RuntimeException e) {
            System.out.print("B");
            throw new RuntimeException();
        } finally {
            System.out.print("C");
        }
    }
}
\end{lstlisting}
	      \textbf{输出：}\underline{\hspace{5cm}}

	\item
	      \begin{lstlisting}[style=JavaExam]
import java.util.*;
public class TestD {
    public static void main(String[] args) {
        Map<String, Integer> map = new HashMap<>();
        map.put("a", 1);
        map.put("b", 2);
        map.put("c", 3);
        map.computeIfPresent("b", (k, v) -> v + 10);
        map.computeIfAbsent("d", k -> 4);
        System.out.println(map);
    }
}
\end{lstlisting}
	      \textbf{输出：}\underline{\hspace{5cm}}

	\item
	      \begin{lstlisting}[style=JavaExam]
public class TestE {
    static int x = 10;
    static { x += 5; }
    static { x /= 3; }
    public static void main(String[] args) {
        System.out.println(x);
    }
}
\end{lstlisting}
	      \textbf{输出：}\underline{\hspace{5cm}}
\end{enumerate}

\newpage

\section*{五、编程题（共 2 小题，每小题 20 分，共 40 分）}
\textit{（要求代码完整、结构清晰、注释合理）}

\begin{enumerate}[label=\arabic*., itemsep=1em]
	\item \textbf{设计一个简单的考试系统} \\
	      定义抽象类 `Question`，包含属性 `String content`（题目内容）和 `double score`（分值）。提供抽象方法 `boolean checkAnswer(String answer)` 用于判断答案是否正确。\\
	      派生两个子类：
	      \begin{itemize}
		      \item `ChoiceQuestion`（选择题）：包含属性 `String[] options`（选项数组）和 `String correctAnswer`（正确答案，如 "A"）。实现 `checkAnswer` 方法。
		      \item `TrueFalseQuestion`（判断题）：包含属性 `boolean correctAnswer`。实现 `checkAnswer` 方法（传入 "true" 或 "false" 字符串）。
	      \end{itemize}
	      定义 `Exam` 类，使用 `ArrayList<Question>` 存储题目，提供方法：
	      \begin{itemize}
		      \item `void addQuestion(Question q)`
		      \item `void takeExam()`：遍历题目，显示题目内容（选择题显示选项），从控制台读取用户答案，计算总分并输出。
	      \end{itemize}
	      在 `Main` 类中测试，创建至少 3 道选择题和 2 道判断题，模拟考试过程。
	      \vspace{6cm}

	\item \textbf{实现一个简单的电影管理系统} \\
	      定义 `Movie` 类，属性包括 `String id`（编号）、`String title`（片名）、`String director`（导演）、`int year`（上映年份）、`double rating`（评分）。提供构造方法、getter/setter，并重写 `toString()` 方法。\\
	      定义 `MovieManager` 类，使用 `HashMap<String, Movie>` 存储电影（键为编号）。提供以下功能：
	      \begin{itemize}
		      \item `void addMovie(Movie m)`：添加电影，若编号已存在则抛出异常。
		      \item `void removeMovie(String id)`：删除电影。
		      \item `Movie findMovie(String id)`：根据编号查找。
		      \item `List<Movie> findMoviesByYear(int year)`：根据年份查找电影列表。
		      \item `List<Movie> getTopRatedMovies(int n)`：返回评分最高的 n 部电影（按评分降序）。
		      \item `void displayAll()`：显示所有电影信息。
	      \end{itemize}
	      在测试类中演示所有功能。
	      \vspace{6cm}
\end{enumerate}

\begin{center}
	\textbf{—— 试卷结束，祝考试顺利 ——}
\end{center}

\end{document}
